% !TeX root = ../papers/nc-p3.tex
\section{有限トーラス窓に特殊化する対称性を指定しない全窓}
\label{sec:finite-torus-slices}
定理\ref{thm:simple-weight-finite-windows}の曲面で $(b,c)=(2,3)$、曲線で $b=2$、
定理\ref{thm:double-middle-finite-windows}で $\alpha=2$ の閉軌道を $A,B,C$ とする。
ここで $b,c$ は六単元正規形の座標である。
\begin{theorem}\label{thm:finite-torus-slices}
二隣接切断の階数6軌跡は三点で滑らかで、アフィン次元は $76,75,75$、各点を通る既約成分は一意である。
$A$ の成分と $B$ の成分は異なり、両者とも $(u,v)=(1/2,1/2)$ の軌道を含む。
一般の添字保持類の次元はそれぞれ $15,14,14$ で、各幾何学的一般点は標準正則である。

各閉軌道に特殊化する全窓は有限点グラフを持ち、トーラスの存在を仮定せず全て正則である。
閉軌道自身を除くそれらの添字保持類は、集合として次の通りである。
\begin{enumerate}
\item $A$ では $\mathbb P(1,2,3,3,4,4,5)$ 一つで、全て最小周期4である。
\item $B$ では同じ重み付き射影空間の二つの互いに素なコピーである。
\item $C$ では $\mathbb P(1,1,2,2,2,3,3)$ の対合商である。重み $1,2,3$ ごとに
符号を変える座標数はそれぞれ一つで、他は固定する。
\end{enumerate}
同じ計算は各表示族の稠密開部分で成立する。
\end{theorem}
\begin{proof}[非被約点での余法線と接空間]
補題\ref{lem:finite-flat-conormals}の証明では四文字積の評価の階数が $p$ のとき、
有限平坦な点環上で作る恒等式の係数行列から $p$ 本の独立行を取れる。
従って第一切断の滑らかな156次元軌跡上で第二切断のSchur方程式は局所的に $100-p$ 本以下で生成される。
$p=19$ の場合も同じ単元小行列の消去でよい。合わせた微分の階数が $200-p$ なら滑らかである。

極小点のアフィンチャートの局所環は前二節の消去式で、単重みなら重み $1,2$ の二変数上の重み $4,5$ の二式、
中央二重なら二変数上の二本の三次式である。最高点にも反転した式を用いる。
四つの一文字の積を、順に位置 $3,0,1,2$ の座標を使い、これらの有限環で剰余にする。
$C$ では中間二つの単純点の評価も加える。256列の階数は順に $20,19,19$ となる。
一方、各切断の左右核のベクトル $l,r$ に対する $l^{\mathsf t}(\delta w)r$ の全てで一次条件を作り、
位置別 $\mathfrak{gl}_4$ の軌道方向を除く。最上重みを引き、
正負の同じ絶対重みをそれぞれ一回ずつ数えたときの各側の多重度は次表になる。
\[
\begin{array}{c|c|ccccc}
 &0&\pm1&\pm2&\pm3&\pm4&\pm5\\\hline
 A&2&1&1&2&2&1\\
 B&1&1&1&2&2&1\\
 C&1&2&3&2&0&0
\end{array}
\]
これ以外の重みは零である。アフィン軌道は各60次元なので上記の微分階数が得られる。
正確な剰余・階数は \path{computations/verify_finite_torus_slices.py} による。
各族は既約なので、その全体が標本での一意な成分に入る。稠密開部分では滑らかさが保たれ、
トーラス同変接束の重み階数は一定である。
二族の共通パラメータが両成分に入るのは先の正規形による。
一般点の正則性は命題\ref{prop:unrestricted-good-reduction}の次数ごとの持上げによる。

小次元側の二つの表示は両対称化もできる。位置0,1にそれぞれ作用させて反対相手との交換を対称にする
対角行列は、単重み曲線では $\operatorname{diag}(b,1,1,1),I$、中央二重では
$\operatorname{diag}(1,1,\alpha,1),\operatorname{diag}(-1/\alpha,1-1/\alpha,-1+1/\alpha,1)$ である。
これは全語の係数比較で直接得られる。
\end{proof}

\begin{proof}[閉軌道の全ファイバー]
射影基底群 $G$ の半安定軌跡で各閉点のアフィン不変近傍を使う。
$H$ を閉軌道の安定群とする。標数零の滑らかな点でのLunaのエタールスライス定理により、
軌道に垂直な上表の $H$ 表現 $N$ とスライスの間に強エタールな局所模型がある。
商で閉点の剰余体は $k$ なので、商上のその一点のファイバーを取ると、
実際の $G$ ファイバーは $G\times^H (N\mathbin{\times}_{N/\!/H}\{0\})$ と同型である。
すなわち線形近似だけでなく、この一点に特殊化する全軌道の同型類を与える。
ランク6と有限辺は閉軌道を含む不変開条件であり全ファイバーで保たれる。
正則性は閉軌道への付値特殊化と命題\ref{prop:unrestricted-good-reduction}から従う。

$H^\circ=T$ が一次元トーラスなので、線形零ファイバーの台は正・負の重み空間の和集合である。
零重みは零でなければならず、正負がともに非零なら次数を揃えた積の不変式がある。
有限拡大 $H/T$ を加えても零ファイバーの台は変わらない。
$A,B$ で向きを保つ $H$ は $T$ である。実際四固有線は単純で、先の12主元の指標格子は飽和している。
$A$ は反転があり正負を交換し、$B$ では $b\ne1-b$ なので反転がない。
$C$ で向きを保ち二つの中間線を交換するには $\alpha=\alpha^{-1}$ が必要で、これは許されない。
従って対角正規化の残余だけで、$C_0=I$ を保つ条件が比の二乗を1にする。
$T$ のほかに一方の中間文字だけ四位置で符号反転する対合が生じる。
重みを反転する同型は極の交換と中間二線の交換を組み合わせて得られるため正負は一つになる。
符号を変える接方向の多重度は上記コードでテンソル単位をその文字数の偶奇にも分けて得る。
非零方向をトーラスで割れば記載した重み付き射影空間となる。
$A$ の場合に一添字または二添字の移動で一般軌道が自分に戻れば、その閉極限も同じ移動で自分に戻る。
この閉極限は周期4なのでファイバー内のどの軌道も周期4である。
いずれの法表現にも非零重みがあり、零でない一般方向の安定群は有限であるから
一般軌道のアフィン次元は61、従って表示した一般同型類次元となる。
\end{proof}
